\section{Results}
\label{sec:results}

We organize the results around five claims, each addressed by a single
section-conditional slice (\S\ref{sec:limitations} L7 discusses the
slicing rationale): the headline 2D \emph{Vulnerability Fingerprint}
(\S\ref{sec:results-fingerprint}); its scale invariance on ogbn-arxiv
(\S\ref{sec:results-scaling}); the architectural \emph{Shard Protection
Effect} of the Partition family (\S\ref{sec:results-shard}); the
\emph{hop-distance decay} signature of approximation-error propagation
(\S\ref{sec:results-hop}); and a privacy-axis decoupling via the
\emph{MIA audit} (\S\ref{sec:results-mia}).

\subsection{Vulnerability Fingerprint}
\label{sec:results-fingerprint}

\begin{figure}[t]
  \centering
  \includegraphics[width=0.78\linewidth]{figures/FIG-3_Spectrum.pdf}
  \caption{\textbf{Vulnerability Fingerprint} on the canonical cell
  (Cora/GCN/$r{=}0.05$, 5 seeds). Each point is a method-level cluster
  centre with a 1$\sigma$ ellipse over seeds; axes are the
  attack-specific F1 effects $\Delta F^{\text{attack}}_{\mathrm{IM}}$
  (structural-spread, $\alpha{=}0$) and
  $\Delta F^{\text{attack}}_{\mathrm{TracIn}}$ (gradient-signal,
  $\alpha{=}1$). Marker shape distinguishes pair members within each
  OpenGU category; chords visualise within-pair distance.
  \interim{Refresh with Phase~B numbers; current placeholder uses
  pre-Phase-B coordinates.}}
  \label{fig:fingerprint}
\end{figure}

The dominant predictor of vulnerability is the \emph{mechanism family},
modulated by within-family mechanism choice. We project each cell
onto the IM/TracIn axes, yielding a 2D coordinate
$(\Delta F^{\text{attack}}_{\mathrm{IM}},\,
\Delta F^{\text{attack}}_{\mathrm{TracIn}})$ where each component is a
paired effect (strategy F1 minus same-seed random F1; see
\S\ref{sec:diagnostics}). The axes are basis directions:
$\alpha{=}0$ and $\alpha{=}1$ endpoints of the Hybrid score
(Eq.~\ref{eq:hybrid}) project onto the IM and TracIn axes respectively.
The architectural F1-shift term $\Delta F_{\text{arch}}$ is reported
separately in \S\ref{sec:results-shard}.

\paragraph{Reading the axes as access tiers.}
The two axes have distinct realistic-attacker readings
(Table~\ref{tab:access}). The IM axis is what an \textbf{L1
structural-only} attacker---a coalition with the public graph but no
model access---can achieve at zero additional cost. The TracIn axis is
the \textbf{L2-direct} upper bound (production weights known); the
realistic L2-\textsc{surrogate} attacker pays one shadow-model training
run for some fraction of this number, with the gap set by
cross-architecture transferability. We therefore read the fingerprint
as a \emph{floor-and-ceiling} statement: \emph{any} attacker, however
weak, can reach the IM-axis number; \emph{no} L2 attacker can exceed
the TracIn-axis number; an honest realistic threat sits between the two
with a transferability discount.

\paragraph{Three-stage clustering.}
We test three nested claims of decreasing strength.
\textbf{Stage~1 (method coherence)}: each method's 5-seed cloud forms a
tight cluster; selector noise is bounded after the IM-seed decoupling
fix (\S\ref{app:imseed}).
\textbf{Stage~2a (intra-family coherence)}: within each OpenGU category,
do paired methods cluster? We predict
\emph{convergent} for the Partition pair (GraphEraser, GraphRevoker)
and \emph{outlier} for the IF and Learning pairs.
\textbf{Stage~2b (outlier distance)}: distance(MEGU, GNNDelete) and
distance(IDEA, GIF) in fingerprint space quantifies how far the
within-family alternative decouples from the family centroid.
The contrast across the three pair signatures is itself a headline
insight: \emph{OpenGU categories are not equally homogeneous}.

\paragraph{What the fingerprint shows.}
Figure~\ref{fig:fingerprint} reads as follows.
\textbf{GNNDelete} sits in Q1, with both axes large
(\interim{$\Delta F^{\mathrm{attack}}_{\mathrm{IM}}{=}+16.9$},
\interim{$\Delta F^{\mathrm{attack}}_{\mathrm{TracIn}}{=}+X.X$};
$p$\interim{$<$0.001}, $N{=}5$): the loss surface is directly
perturbed by both gradient-signal and structural-spread selectors.
\textbf{GIF} sits along the TracIn axis with a small but significant
TracIn projection (\interim{$+4.17\pm 0.76$}, $p$\interim{$=$0.0001})
and near-null IM projection: the Newton step responds to gradient
signal but not to diffusion spread.
\textbf{GraphEraser} and \textbf{GraphRevoker} cluster near the origin
with non-positive coordinates---no above-baseline attack signal---a
geometric statement of Shard Protection (\S\ref{sec:results-shard}).
\textbf{MEGU} and \textbf{IDEA}, despite belonging to the
otherwise-vulnerable Learning and IF families, sit close to the origin
($<\interim{1\%}$ on both axes), with cluster centres
\interim{$\sim 5\times$} farther from their family centroids than the
canonical members. This is the within-family heterogeneity finding.
A 1D projection of the same data (the prior ``vulnerability spectrum''
view) is included as Appendix~\ref{app:results-1d}.

\subsection{Scaling Validation on ogbn-arxiv}
\label{sec:results-scaling}

\begin{figure}[t]
  \centering
  \includegraphics[width=0.84\linewidth]{figures/FIG-1_Generalization.pdf}
  \caption{\textbf{Scale invariance of the fingerprint geometry.}
  Same 2D coordinate per method across Cora, Citeseer, and ogbn-arxiv
  (GCN backbone, $r{=}0.05$). Family clusters retain their relative
  ordering; outlier distance for MEGU/IDEA persists at 169K nodes.
  \interim{Phase~B.2 numbers pending; arxiv markers are placeholders.}}
  \label{fig:scaling}
\end{figure}

We replicate the canonical-cell fingerprint on ogbn-arxiv with 3 seeds
(\S\ref{sec:setup-arxiv}). The family ordering and the relative
positions of intra-family pairs are preserved at the
$60\times$-larger scale (Figure~\ref{fig:scaling}): GNNDelete remains
in Q1; GIF retains a TracIn-axis lean; the partition pair stays near
the origin or in the protective region; MEGU and IDEA remain
near-origin outliers. We pass the 11-item feasibility gate of
\interim{Phase~B.1} (definitions in
\S\ref{app:results}). Per-cell tables, including paired-effect CI
widths under the reduced seed count, appear in
Appendix~\ref{app:results}.

\subsection{The Shard Protection Effect}
\label{sec:results-shard}

\begin{figure}[t]
  \centering
  \includegraphics[width=0.92\linewidth]{figures/FIG-2_Scaling.pdf}
  \caption{\textbf{Architectural F1 shift $\Delta F_{\text{arch}}$ per
  method} (Cora/GCN, $r{=}0.05$). Bars below zero mark Shard
  Protection: the Partition pair's downstream F1 \emph{increases} when
  any small random set is unlearned. The bar value is the paired-effect
  reference floor; the fingerprint of \S\ref{sec:results-fingerprint}
  reports the orthogonal $\Delta F_{\text{attack}}$ term.
  \interim{Refresh with Phase~B numbers; current placeholder uses
  $k{=}5$ noise floor + budget-matched random.}}
  \label{fig:shard-protection}
\end{figure}

We separate the \emph{architectural} F1-shift term
$\Delta F_{\text{arch}}$ (Eq.~\ref{eq:fshift}) from the attack-specific
fingerprint of \S\ref{sec:results-fingerprint}. For the Partition
family, $\Delta F_{\text{arch}}<0$: random deletion at $r=0.05$
\emph{improves} F1 by \interim{$3$--$6\%$}. We name this signature the
\textbf{Shard Protection Effect}.

\paragraph{What it is and what it is not.}
Shard Protection is \emph{not} an attack-specific finding---our attack
does not cause it and does not amplify it. It is a property of running
a partition-aggregate model on $\mathcal{D}_{\text{tr}}\setminus S$ for
\emph{any} reasonable $S$ at small budget, already visible in the
$k{=}5$ noise-floor baseline (\interim{$+0.0$ vs $+X.X$}). The
fingerprint cleanly separates the two:
$\Delta F_{\text{attack}}\approx 0$ on shard methods (origin in
Figure~\ref{fig:fingerprint}), while $\Delta F_{\text{arch}}<0$
(Figure~\ref{fig:shard-protection}). Reporting both terms is what makes
the honesty of the claim visible.

\paragraph{Mechanism.}
GraphEraser partitions the graph into $C$ disjoint shards, trains a
sub-classifier per shard, and aggregates by learnable attention.
Deleting a training node only affects the shard it belongs to; the
aggregator's other shards are untouched. Empirically, deletion of
high-degree (heuristic) or high-influence (TracIn/IM) nodes
\emph{purifies} a shard---removing nodes that span class boundaries
leads to cleaner sub-classifiers, and the attention aggregator is
robust enough to absorb the change. The attacker's upper bound is
therefore set by random selection, which uniformly thins each shard.
GraphRevoker exhibits the same pattern within
\interim{$\le 0.5\%$} F1 (Appendix~\ref{app:shard}); the protection is
mechanism-driven rather than implementation-specific. Robustness to
shard-count $C\in\{5,10,20\}$ is reported in Appendix~\ref{app:shard}.

\paragraph{Implication.}
For shard-based GU, adversarial deletion is \emph{not} a meaningful
threat at small budgets---and may even be a useful regularizer. We do
not interpret this as a recommendation to deploy partition-aggregate GU
unconditionally; the trade-off sits in privacy guarantees rather than
utility. A defender's reading is in \S\ref{sec:broader}.

\subsection{Collateral Damage and Hop-Distance Decay}
\label{sec:results-hop}

\begin{figure}[t]
  \centering
  \includegraphics[width=0.48\linewidth]{figures/FIG-4a_Significance.pdf}
  \hfill
  \includegraphics[width=0.48\linewidth]{figures/FIG-4b_Effect.pdf}
  \caption{\textbf{Effect-size and significance heatmap} (5-seed paired
  $t$-test, Cora/GCN, $r{=}0.05$). Each cell is colored by the $p$-value
  of the one-sided test $\mu\le 0$. Diagonal patterns confirm the
  fingerprint geometry of \S\ref{sec:results-fingerprint} and motivate
  the hop-decay reading below. \interim{Refresh with Phase~B
  (decoupled IM seed) numbers; pre-fix data had Jaccard $=0.13$ between
  IM seeds.}}
  \label{fig:significance}
\end{figure}

We report the three orthogonal collateral metrics introduced in
\S\ref{sec:diagnostics}.

\paragraph{Significance.}
Figure~\ref{fig:significance} reports the paired $t$-test over five
seeds for every (method, strategy) cell. With the
\texttt{im\_selector\_seed} fix (\S\ref{app:imseed}), all but
\interim{one} of the GNNDelete IM/Hybrid cells clear $p<0.05$. The GIF
row clears $p<0.001$ for TracIn but Hybrid is borderline, consistent
with the Newton-step ceiling reading of
\S\ref{sec:results-fingerprint}. The MEGU and IDEA rows show no cell
clearing $p<0.05$ across any strategy.

\paragraph{Retrain gap.}
Across methods, the F1 drop is \emph{not} explained by simply having
removed informative nodes: retraining from scratch on
$\mathcal{D}_{\mathrm{tr}}\setminus S$ recovers baseline-level F1
(retrain gap \interim{$<$2\%} for GIF and the partition pair). For
GNNDelete the retrain gap is \interim{20.5\%} on average, confirming
that the F1 collapse comes from the unlearning algorithm's
approximation error rather than from data loss.

\paragraph{Hop-distance decay.}
We report flip rate vs.\ retrain (Eq.~\ref{eq:hopflip}) in three
RF-relative buckets per backbone (within-RF, boundary, outside-RF). On
Cora/GCN ($L=2$), GIF and the partition pair exhibit \emph{sharp
decay}: flip rate concentrates within-RF and falls to near zero
outside-RF---a \emph{localized} signature. GNNDelete exhibits
\emph{near-uniform} flip rate across all three buckets---a
\emph{propagating} signature, indicating that its representation-level
objective drives approximation drift beyond the receptive field. MEGU
and IDEA stay near zero in every bucket, consistent with their
fingerprint position. The same pattern survives the backbone change to
ogbn-arxiv ($L=3$), where the buckets shift in absolute hop count but
preserve the within-RF / boundary / outside-RF semantics. Per-family
decay curves with seed-level error bars and the absolute 4-bucket
breakdown appear in Appendix~\ref{app:results}.

\subsection{MIA Audit: Privacy and Utility Decouple}
\label{sec:results-mia}

\begin{table}[t]
\centering
\small
\caption{\textbf{MIA AUC vs.\ paired F1 effect.} Privacy axis (MIA AUC
post-unlearning, shadow-model
attack~\cite{olatunji2021membership}) and utility axis
($\Delta F^{\text{attack}}$) decouple. Cora/GCN, $r{=}0.05$, 5 seeds,
post bug-fix. \interim{Refresh with Phase~B values.}}
\label{tab:mia}
\begin{tabular}{@{}lllrr@{}}
\toprule
Family & Method & MIA AUC & $\Delta F^{\text{attack}}_{\mathrm{IM}}$ (\%) & Reading \\
\midrule
Partition & GraphEraser  & \interim{0.60} & \interim{$+0.0$}     & non-private, robust utility \\
Partition & GraphRevoker & \interim{0.6X} & \interim{$+0.0$}     & non-private, robust utility \\
IF        & GIF          & \interim{0.5X} & \interim{$+X.X$}     & weak privacy, weak attack \\
IF        & IDEA         & \interim{0.55} & \interim{$+0.X$}     & weak privacy, null attack \\
Learning  & GNNDelete    & \interim{0.6X} & \interim{$+16.9$}    & weak privacy, severe attack \\
Learning  & MEGU         & \interim{0.37} & \interim{$+0.X$}     & strong privacy, null attack \\
\bottomrule
\end{tabular}
\end{table}

The privacy axis (MIA AUC after unlearning) and the utility axis
($\Delta F^{\text{attack}}$) are not aligned (Table~\ref{tab:mia}).
MEGU resists the F1 attack while \emph{also} reducing membership
leakage (AUC \interim{$\approx 0.37$}); the partition pair resists the
F1 attack but retains weak-to-moderate membership leakage
(AUC \interim{$\approx 0.60$}); GNNDelete is bad on both axes. This
decoupling supports the framing of MIA as a complementary diagnostic,
not a redundant one, and aligns with the within-family heterogeneity
reading of \S\ref{sec:results-fingerprint}: MEGU and IDEA are not
``robust'' in a single sense---they decouple from \emph{both} the
selector signals our attacks exploit \emph{and} from the baseline
membership-leakage signal of their family-mates.
